% Снимок текущего состояния физики и математики воксельной модели айсберга (Stage 12)
% Компиляция: xelatex (или lualatex); шрифты TeX Gyre
\documentclass[12pt,a4paper]{article}
\usepackage{fontspec}
\usepackage{polyglossia}
\setdefaultlanguage{russian}
\setotherlanguage{english}
\IfFontExistsTF{Times New Roman}{\setmainfont{Times New Roman}}{%
\IfFontExistsTF{Liberation Serif}{\setmainfont{Liberation Serif}}{%
\IfFontExistsTF{STIX Two Text}{\setmainfont{STIX Two Text}}{%
\IfFontExistsTF{Libertinus Serif}{\setmainfont{Libertinus Serif}}{\setmainfont{FreeSerif}}}}}
\IfFontExistsTF{Consolas}{\setmonofont{Consolas}}{%
\IfFontExistsTF{Courier New}{\setmonofont{Courier New}}{%
\IfFontExistsTF{Liberation Mono}{\setmonofont{Liberation Mono}}{%
\IfFontExistsTF{DejaVu Sans Mono}{\setmonofont{DejaVu Sans Mono}}{\setmonofont{FreeMono}}}}}
\usepackage{amsmath,amssymb}
\usepackage{geometry}
\geometry{left=30mm,right=15mm,top=20mm,bottom=20mm}
\usepackage{setspace}
\onehalfspacing
\usepackage{indentfirst}
\setlength{\parindent}{1.25cm}
\usepackage{booktabs,longtable,array}
\usepackage{titlesec}
\titleformat{\section}{\normalfont\fontsize{14}{17}\bfseries}{\thesection}{0.6em}{}
\titleformat{\subsection}{\normalfont\fontsize{12}{15}\bfseries}{\thesubsection}{0.6em}{}
\usepackage[hidelinks]{hyperref}

\newcommand{\Lf}{L_{\mathrm f}}
\newcommand{\rhoi}{\rho_{\mathrm i}}
\newcommand{\rhow}{\rho_{\mathrm w}}
\newcommand{\ci}{c_{\mathrm i}}
\newcommand{\ki}{k_{\mathrm i}}
\newcommand{\Tf}{T_{\mathrm f}}
\newcommand{\code}[1]{\texttt{\small #1}}
\setlength{\LTpre}{6pt}\setlength{\LTpost}{6pt}

\emergencystretch=3em
\begin{document}

\begin{center}
{\fontsize{16}{20}\selectfont\bfseries Воксельная термодинамическая модель айсберга:\\ физическая и математическая постановка в текущей реализации (Stage~12)}\\[1ex]
Состояние на 19 ноября 2026 г. (корректирующий аудит Stage 12.22)
\end{center}

\tableofcontents
\newpage

%=====================================================================
\section{Назначение и границы документа}

Документ фиксирует действующую постановку воксельной модели айсберга (Stage~12): уравнения, разностные схемы, пространственное и временное разрешение, порядок вычислений на шаге, учтённые и неучтённые процессы, места, где физика заменена упрощением (заглушкой), и способы, которыми каждое упрощение может быть снято. Описана реализация в каталоге \code{src/} и приложении \code{app/iceberg\_run.f90} (базовый коммит физики/Fortran: \code{cd33f0f}, Stage 12.20; после него Fortran не менялся). Валидационные/документационные коммиты: \code{18566e9} (Stage 12.22), \code{d7476bd} (Stage 12.22a), \code{34e1f89} (Stage 12.22b); ни один из них не менял Fortran-физику, константы или численные схемы. Текущий HEAD репозитория: \code{34e1f89}. Настоящая сборка 12.22c не меняет физику, константы или схемы; расчёты Stage~12 в Баренцевом море завершены; призматическая модель Stage~11 заморожена и привлекается только как эталон проверок.

Модель односторонне связана с реанализами: океан и атмосфера заданы извне, обратного влияния айсберга (талая вода, охлаждение) на них нет. Разделы 2--12 описывают то, что реально считается; разделы 13--16 --- волны, проверки, ограничения и развитие. Всё, что в разделе \ref{sec:future} названо будущим развитием, в коде не реализовано.

%=====================================================================
\section{Общая структура и порядок шага}\label{sec:order}

Состояние айсберга: поля $\varphi_i$ (объёмная доля льда) и $T_i$ (температура, \textdegree C) на воксельной сетке; положение $(\phi,\lambda)$ и скорость $\mathbf u$. Шаг по времени $\Delta t=3600$~с одинаков для динамики и термодинамики. Вокселизованное состояние является основным; интегральные величины (масса, осадка, надводный борт, площади) вычисляются по нему. Если атмосферный файл не задан, приложение использует постоянную атмосферу; в воксельном шаге автономный воздушный баланс передаётся через \code{autonomous\_air=.true.} (ERA5-поля либо постоянные значения).

Один шаг состоит из следующих операций.
\begin{enumerate}
\item Интерполяция внешних полей в точку айсберга: профили $T_w(z),S_w(z),u(z),v(z)$ из GLORYS12 и приземные поля ERA5 (раздел~\ref{sec:forcing}); расчёт $S_\downarrow$ и $L_\downarrow$ (раздел~\ref{sec:air}).
\item Гидростатика: ватерлиния $z_{wl}$, осадка, надводный борт; маски погружённых и надводных вокселей; классификация открытых граней (разделы~\ref{sec:hydro}, \ref{sec:open}).
\item Для каждой открытой подводной грани --- решение трёхуравненной задачи на границе лёд--вода по температуре вокселя $T^n$: получаются $m,\;T_B,\;S_B$ (раздел~\ref{sec:3eq}); объём таяния грани накапливается в $\Delta\varphi^{\mathrm{bnd}}_i=\sum_f m_f\Delta t/h$.
\item Тепловая задача Кранка--Николсон по всей сетке с граничными условиями: условие Дирихле $T_B$ на подводных гранях, линеаризованное условие Робина на надводных (раздел~\ref{sec:cn}).
\item Фазовое ограничение: воксели с $T>0$ \textdegree C теряют долю льда, избыток энтальпии идёт в таяние, $T\leftarrow0$ (раздел~\ref{sec:melt}).
\item Граничное таяние: $\varphi_i\leftarrow\varphi_i-\min(\Delta\varphi^{\mathrm{bnd}}_i,\,0{,}2\varphi_i)$ (раздел~\ref{sec:melt}).
\item Слив вокселей с $\varphi<\varphi_{\min}$ (передача явного тепла соседям); каждые 1~ч по умолчанию (переопределение \code{VOXEL\_CONNECTIVITY\_HOURS}, максимум 6~ч; в этапе 12.20 использовано 6~ч) --- анализ связности и удаление отколовшихся компонент; вертикальное уплотнение (раздел~\ref{sec:flush}).
\item Динамика: обновление $\mathbf u$ по новой геометрии (площади киля и паруса из вокселей), перемещение айсберга (раздел~\ref{sec:dyn}).
\item Бухгалтерия энергии и массы (раздел~\ref{sec:ledger}).
\end{enumerate}

%=====================================================================
\section{Пространственная и временная дискретизация}\label{sec:grid}

\subsection{Сетка}
Используется равномерная декартова сетка с кубическими вокселями. \textbf{Программные значения по умолчанию} --- $N_x=N_y=N_z=10$, $h=10$~м; \textbf{научная конфигурация} --- $N_x=N_y=N_z=100$, $h=1$~м (то есть $10^6$ ячеек, начальное тело --- куб $100\times100\times100$~м, $\varphi=1$ во всех ячейках, начальная температура $T_0=-10$~\textdegree C). Для ретроспективного расчёта Enderlin23 использована тестовая сетка $8^3$ с $h=12.5$~м. Сетка задаётся через \code{VOXEL\_NX/NY/NZ} и \code{VOXEL\_DX/DY/DZ}. Оси: $x$ и $y$ --- горизонтальные, $z$ --- вверх от дна тела (индекс $k$ растёт вверх). Центры ячеек $x_i=(i-\tfrac12)h$. Объём $V=h^3$; площадь грани $A=h^2$; расстояние по нормали $h_n=h$.

Сетка привязана к телу (не к географической): на нём нет вращения, геометрия меняется только удалением льда. Допускается другая размерность сетки и шаг (параметры запуска); для проверки сходимости использовались $8^3$ и $100^3$.

\subsection{Масштабы и обоснование разрешения}
Температуропроводность льда $\alpha=\ki/(\rhoi\ci)=2{,}2/(910\cdot2100)\approx1{,}15\cdot10^{-6}$~м\textsuperscript2/с. Глубина проникновения тепловой волны за время $t$ порядка $\sqrt{\alpha t}$: за сутки $\approx0{,}3$~м, за 30 суток $\approx1{,}7$~м, за год $\approx6$~м. Шаг $h=1$~м поэтому разрешает сезонный и месячный тепловой пограничный слой, но не разрешает суточный ход температуры поверхности (его глубина меньше $h$). Это принято как осознанное упрощение: в поверхностный воксель вся радиация поступает как поток, суточный ход усредняется (атмосферные поля суточные).

Явная схема устойчива при $\Delta t\le h^2/(6\alpha)\approx1{,}4\cdot10^5$~с для кубической ячейки, то есть при $\Delta t=3600$~с запас на два порядка; ограничением служит не устойчивость, а необходимость связного неявного учёта граничных потоков и фазового ограничения.

\subsection{Время}
$\Delta t=3600$~с для всех подсистем. Внешние поля заданы суточными значениями и интерполируются линейно по времени. Схема по времени для температуры --- Кранк--Николсон (раздел~\ref{sec:cn}); для динамики --- полунеявная (раздел~\ref{sec:dyn}); для геометрии --- явная (потеря льда за шаг ограничена долей $0{,}2$ от $\varphi$ вокселя).

\subsection{Привязка по пространству к внешним данным}
GLORYS12 (горизонтально $1/12^\circ$, 50 уровней) интерполируется билинейно по горизонтали в точку айсберга и линейно по глубине; профиль по глубине $z$ берётся в центре грани относительно ватерлинии $z_{wl}$: глубина грани $d=\max(0,\,z_{wl}-z_f)$, где $z_f$ --- вертикальная координата центра боковой грани, $k h$ для верхней и $(k-1)h$ для нижней грани. Ниже дна профиль экстраполируется последним уровнем со скоростью $u=v=0$.

%=====================================================================
\section{Состояние, гидростатика и связность}\label{sec:hydro}

\subsection{Переменные}
$\varphi_i\in[0,1]$ --- объёмная доля льда в вокселе $i$; воксель активен при $\varphi_i>\varepsilon_\varphi=10^{-12}$ и считается телом при $\varphi_i\ge\varphi_{\min}=0{,}01$. Теплоёмкость вокселя
\begin{equation}
C_i=\varphi_i\,\rhoi\,\ci\,V,
\end{equation}
где $\rhoi=910$~кг/м\textsuperscript3, $\ci=2100$~Дж/(кг$\cdot$К). Энтальпия принимается равной $\ci T$ при $T\le0$ \textdegree C, скрытая теплота $\Lf=3{,}34\cdot10^5$~Дж/кг учитывается при таянии.

\subsection{Гидростатика}
Масса главного тела $M=\rhoi V\sum_i\varphi_i$ (суммирование по ячейкам главной компоненты). Из закона Архимеда объём вытесненной воды $V_{\mathrm{sub}}=M/\rhow$, $\rhow=1028$~кг/м\textsuperscript3. Ватерлиния $z_{wl}$ определяется как уровень, при котором объём твёрдых ячеек ниже него равен $V_{\mathrm{sub}}$ (ячейки считаются твёрдыми целиком). Осадка и надводный борт:
\begin{equation}
D=z_{wl}-z_{\mathrm{bot}},\qquad F_b=z_{\mathrm{top}}-z_{wl},
\end{equation}
где $z_{\mathrm{bot}}=(k_{\mathrm{bot}}-1)h$ и $z_{\mathrm{top}}=k_{\mathrm{top}}h$ --- нижняя и верхняя границы непустых ячеек главной компоненты. Воксель считается погружённым, если $(k-\tfrac12)h<z_{wl}$, иначе надводным.

\subsection{Открытые грани}\label{sec:open}
Грань вокселя открыта, если соседняя ячейка пуста (или это внешняя граница сетки) и эта пустота связана с внешней областью через пустые ячейки (заливка от границы сетки). Изолированные полости (замкнутые пустоты внутри тела) к воде не открыты и потоков не получают.

\subsection{Связность и откол}\label{sec:flush}
Компоненты вокселей с $\varphi>0{,}3$ (\code{VOXEL\_PHI\_CONN}) проверяются каждые 1~ч по умолчанию (\code{VOXEL\_CONNECTIVITY\_HOURS}; максимум 6~ч; в этапе 12.20 использовано 6~ч). Крупнейшая считается главным телом, остальные удаляются как отколовшиеся: их масса и сенсибельная энергия заносятся в бухгалтерию, судьба фрагментов не прослеживается. Воксели с $\varphi<\varphi_{\min}$ сливаются: масса и энергия фиксируются в счётчиках слива, а явное тепло $\sum\rhoi\ci\varphi V T$ распределяется поровну между активными шестью соседями. После слива и после вертикального уплотнения (смещение содержимого вниз на пустые нижние слои, чтобы тело оставалось у нижней границы сетки) включается процедура Раннахера (раздел~\ref{sec:cn}).

%=====================================================================
\section{Динамика дрейфа}\label{sec:dyn}

Уравнение движения (масса $M$ с учётом текущей геометрии, горизонтальная скорость $\mathbf u$):
\begin{equation}
M\frac{d\mathbf u}{dt}=\mathbf F_a+\mathbf F_w+\mathbf F_c+\mathbf F_p .
\end{equation}

\begin{itemize}
\item Ветер: $\mathbf F_a=\tfrac12\rho_a\left(C_a A_{\mathrm{sail}}+C^{s}_aA_{\mathrm{top}}\right)|\mathbf V_{10}-\mathbf u|(\mathbf V_{10}-\mathbf u)$, где $\rho_a=1{,}225$~кг/м\textsuperscript3, $C_a=1{,}3$ (форменное сопротивление парусной части), $C^s_a=2{,}2\cdot10^{-3}$ (кожное сопротивление верхней поверхности), $A_{\mathrm{sail}}$ --- вертикальная площадь надводной части, $A_{\mathrm{top}}$ --- горизонтальная площадь верха.
\item Вода (по слоям): $\mathbf F_w=\sum_k\tfrac12\rhow\,C_w\,A^{(k)}_{\mathrm{keel}}\,|\mathbf u_k-\mathbf u|(\mathbf u_k-\mathbf u)+\tfrac12\rhow C^s_wA_{\mathrm{bot}}|\mathbf u_{\mathrm b}-\mathbf u|(\mathbf u_{\mathrm b}-\mathbf u)$, где $C_w=0{,}9$, $C^s_w=5{,}5\cdot10^{-3}$, $A^{(k)}_{\mathrm{keel}}$ --- вертикальная площадь вокселей слоя $k$, $\mathbf u_k$ --- течение на глубине слоя, $\mathbf u_{\mathrm b}$ --- течение у дна тела, $A_{\mathrm{bot}}$ --- горизонтальная нижняя площадь.
\item Кориолис: $\mathbf F_c=-Mf\,\mathbf k\times\mathbf u$, $f=2\Omega\sin\phi$.
\item Градиент давления $\mathbf F_p$ в режиме реанализа не задаётся (нулевой): геострофическая часть уже входит в течения GLORYS12.
\end{itemize}
Выбор $C_a$, $C_w$ и кожных коэффициентов следует \cite{bigg1997,keghouche2009}; послойный учёт сопротивления --- по подходу NEMO-ICB \cite{marsh2015}. Сопротивление и Кориолис интегрируются полунеявно: скорость нового шага находится из линейной системы $2\times2$, в которую сопротивление входит в линеаризованном виде (скорость относительного течения берётся с предыдущего шага). Положение обновляется явно:
\begin{equation}
\Delta\phi=\frac{v\Delta t}{R},\qquad\Delta\lambda=\frac{u\Delta t}{R\cos\phi},
\end{equation}
где $R$ --- радиус Земли.

\textbf{Заглушки динамики:} волновая сила и стоксов дрейф выключены; сопротивление морского льда отсутствует; опрокидывание и вращение тела не моделируются (вычисляется только метацентрическая оценка $GM=BM-BG$, $BM=I_{wp}/V_{\mathrm{sub}}$, как диагностика); дрейф не валидирован по наблюдениям.

%=====================================================================
\section{Внешние данные}\label{sec:forcing}

\begin{itemize}
\item \textbf{Океан:} GLORYS12, суточные поля $1/12^\circ$: $T_w,S_w,u,v$ на 50 уровнях.
\item \textbf{Атмосфера:} ERA5, суточные средние: $u_{10},v_{10}$, $t_{2m}$, $d_{2m}$, общая облачность $tcc$, давление $msl$, снегопад. Потоки радиации ERA5 (SSRD, STRD) не используются.
\item \textbf{Рельеф дна:} IBCAO V5.2 --- для проверки грунтования.
\end{itemize}
Связь односторонняя, оффлайн: айсберг не изменяет поля океана и атмосферы.

%=====================================================================
\section{Тепловая задача: уравнение и схема Кранка--Николсон}\label{sec:cn}

\subsection{Полудискретное уравнение}
Для льда без фазового перехода справедливо уравнение теплопроводности $\rhoi\ci\,\partial T/\partial t=\nabla\!\cdot(\ki\nabla T)$. Для воксельной модели оно записано в форме баланса по ячейке (метод конечных объёмов на кубической сетке, центральные разности второго порядка):
\begin{equation}
C_i\frac{dT_i}{dt}=\sum_{j\in\mathcal N(i)}G_{ij}\,(T_j-T_i)+\sum_{f\in\mathcal B(i)}Q_{i,f},\qquad
G_{ij}=\frac{\ki A}{h_n}\,\min(\varphi_i,\varphi_j),
\end{equation}
где $\mathcal N(i)$ --- шесть соседей-ячеек со льдом ($\varphi_j>\varepsilon_\varphi$), $G_{ij}$ --- проводимость между ячейками, $\mathcal B(i)$ --- открытые грани ячейки, $Q_{i,f}$ --- тепловой поток через граничную грань, Вт; $\ki=2{,}2$~Вт/(м$\cdot$К). Коэффициент $\min(\varphi_i,\varphi_j)$ ослабляет проводимость через частично заполненные ячейки; для $\varphi\equiv1$ уравнение переходит в стандартную семиточечную аппроксимацию лапласиана $\alpha\nabla^2T$ (источник схемы: \cite{samarskii1977}).

\subsection{Схема по времени}
Применяется $\theta$-схема (при $\theta=0{,}5$ --- Кранк--Николсон \cite{crank1947}):
\begin{equation}
\frac{C_i}{\Delta t}\left(T_i^{n+1}-T_i^{n}\right)=\theta\,F_i\!\left(T^{n+1}\right)+(1-\theta)\,F_i\!\left(T^{n}\right),\qquad\theta=0{,}5,
\end{equation}
где $F_i$ --- правая часть полудискретного уравнения. При $\theta=1$ получается неявная схема Эйлера; при $\theta=0$ --- явная. Схема безусловно устойчива (при $\theta\ge0{,}5$), имеет второй порядок по времени; $C_i$ берётся на начало шага.

\subsection{Линейная система}
После подстановки получается линейная система $\mathbf M\mathbf T^{n+1}=\mathbf r$ с семидиагональной матрицей:
\begin{equation}
M_{ii}=\frac{C_i}{\Delta t}+\theta\sum_{j}G_{ij}+\theta\,G^{\mathrm{sub}}_i-\sum_{f\in\mathrm{air}}A\,\frac{\partial Q}{\partial T_s},\qquad M_{ij}=-\theta\,G_{ij},
\end{equation}
\begin{equation}
\begin{split}r_i={}&\frac{C_i}{\Delta t}T_i^n+(1-\theta)\sum_jG_{ij}\left(T_j^n-T_i^n\right)\\&+\sum_{f\in\mathrm{sub}}G^{\mathrm{b}}_f\left[T_{B,f}-(1-\theta)T_i^n\right]+\sum_{f\in\mathrm{air}}A\left[Q_0-\frac{\partial Q}{\partial T_s}T_i^n\right],\end{split}
\end{equation}
где $G^{\mathrm b}_f=2\ki A/h_n$ --- проводимость полуячейки до границы, $G^{\mathrm{sub}}_i=\sum_fG^{\mathrm b}_f$ по открытым подводным граням, $T_{B,f}$ --- температура границы (раздел~\ref{sec:3eq}), $Q_0$ и $\partial Q/\partial T_s$ --- воздушный поток и его производная при $T_s=T_i^n$ (раздел~\ref{sec:air}). Пустые ячейки заменены тождественными уравнениями $T_i^{n+1}=T_i^n$.

Подводные грани входят с весом $\theta$ (стандартный Кранк--Николсон), воздушные --- в линеаризованном виде (один шаг Ньютона от $T_i^n$), то есть полностью неявно относительно $T^{n+1}$. Матрица симметрична и положительно определена и решается методом сопряжённых градиентов с диагональным (Якоби) предобуславливанием \cite{hestenes1952}; начальное приближение $T^n$; критерий $\|r\|/\|b\|<10^{-10}$ (\code{ICEBERG\_CG\_TOL}); по умолчанию максимум $10\,N_xN_yN_z$ итераций (для $100^3$ это $10^7$), переопределяется через \code{ICEBERG\_CG\_MAXIT}; фактический solver использует не менее 10 итераций.

\subsection{Процедура Раннахера}
Схема Кранка--Николсон при нерегулярных данных даёт слабо затухающие осцилляции. После любого события, меняющего геометрию (слив вокселя, уплотнение), на следующие два шага вместо $\theta=0{,}5$ используется $\theta=1$ \cite{rannacher1984}.

\subsection{Явная схема}
Для сопоставления сохранена явная схема (\code{VOXEL\_THERMAL\_SCHEME=explicit}): шаг по времени дробится на подшаги $\Delta t_{\mathrm{sub}}\le h^2/(6\alpha)$, на границе применяется поток явным образом. В расчётах Stage~12 по умолчанию используется Кранк--Николсон.

%=====================================================================
\section{Граница с водой}\label{sec:3eq}

\subsection{Три уравнения интерфейса}
На каждой открытой подводной грани решается задача о тепло- и солеобмене на границе лёд--вода \cite{hollandjenkins1999,jenkins2010}. Неизвестные: скорость отступания границы $m\ge0$, м/с, солёность $S_B$ и температура $T_B$ границы.
\begin{align}
&\rhow\,\gamma_S\,(S_w-S_B)=\rhoi\,m\,S_B &&\text{(солевой баланс, } S_{\mathrm{ice}}\approx0),\label{eq:salt}\\
&T_B=\Tf(S_B,P)=\left(a_0+a_1\sqrt{S_B}-a_2S_B\right)S_B+b\,P &&\text{(температура замерзания)},\label{eq:tf}\\
&\rhow c_w\,\gamma_T\,(T_w-T_B)-\frac{2\ki}{h_n}\left(T_B-T_{\mathrm v}\right)=\rhoi\,\Lf\,m &&\text{(тепловой баланс)}.\label{eq:heat}
\end{align}
Здесь $S_w,T_w$ --- солёность и температура воды вдали от стенки на глубине грани; $T_{\mathrm v}$ --- температура приграничного вокселя; $c_w=3974$~Дж/(кг$\cdot$К); $P$ --- давление на глубине грани, дбар, $P=\rhow gd/10^4$; коэффициенты EOS-80 для температуры замерзания: $a_0=-0{,}0575$~\textdegree C/PSU, $a_1=1{,}710523\cdot10^{-3}$, $a_2=2{,}154996\cdot10^{-4}$, $b=-7{,}53\cdot10^{-4}$~\textdegree C/дбар (солёность в формуле --- в PSU, в уравнении (\ref{eq:salt}) --- в массовых долях); $g=9{,}80665$~м/с\textsuperscript2.

Обменные скорости:
\begin{equation}
\gamma_T=\max\!\left(K_T u_{\mathrm{rel}},\;\gamma_T^{\mathrm{free}}\right),\qquad
\gamma_S=\max\!\left(K_S u_{\mathrm{rel}},\;\gamma_S^{\mathrm{free}}\right),
\end{equation}
$K_T=1{,}1\cdot10^{-3}$, $K_S=3{,}1\cdot10^{-5}$ (безразмерные, по \cite{jenkins2010}), $u_{\mathrm{rel}}=|\mathbf u_w(d)-\mathbf u|$ --- скорость воды на глубине грани относительно айсберга; $\gamma_T^{\mathrm{free}}=10^{-5}$~м/с, $\gamma_S^{\mathrm{free}}=10^{-6}$~м/с --- нижние пределы, имитирующие свободную конвекцию при малой скорости $u_{\mathrm{rel}}$ (журнал решений D-62; в коде ссылка на Hellmer и Olbers \cite{hellmer1989}); это физические нижние оценки, а не результат калибровки.

\subsection{Решение}
Из (\ref{eq:salt}) $S_B(m)=\rhow\gamma_S S_w/(\rhow\gamma_S+\rhoi m)$; тогда (\ref{eq:tf}) задаёт $T_B(m)$, и (\ref{eq:heat}) превращается в $F(m)=0$ с монотонно убывающей $F$. Решение находится бисекцией (до 80 итераций, точность по $m$ порядка $10^{-13}$~м/с). Если $F(0)\le0$, то $m=0$: намерзание на границе не моделируется.

\subsection{Использование результатов}
Для каждой грани сохраняются $T_B,S_B$, которые идут в тепловую систему как условие Дирихле (раздел~\ref{sec:cn}). Скорость отступания переводится в долю льда: $\Delta\varphi=m\Delta t/h_n$. Суммарные скорости таяния: базальная (грани $d=5,6$) и боковая (остальные), в кг/с, $\rhoi A m$.

%=====================================================================
\section{Граница с воздухом}\label{sec:air}

Надводные грани получают тепловой поток в виде автономного линеаризованного баланса (поток положителен, если направлен в лёд):
\begin{equation}
Q=Q_{SW}+Q_{LW}+Q_{SH}+Q_{LH},\qquad Q(T_s)\approx Q_0+\frac{\partial Q}{\partial T_s}\,(T_s-T_i^n).
\end{equation}

\paragraph{Коротковолновая радиация.}
\begin{equation}
Q_{SW}=(1-\alpha)\,\chi\,S_\downarrow,\qquad
S_\downarrow=S_0\cos Z\;\mathcal T_R\,\mathcal T_{wv}\,\mathcal T_a\,\mathcal T_c .
\end{equation}
$S_0=1353$~Вт/м\textsuperscript2 --- солнечная постоянная; $Z$ --- зенитный угол Солнца (по формулам Спенсера \cite{spencer1971}: склонение и уравнение времени); $\alpha=0{,}7$ --- альбедо; $\chi$ --- геометрический фактор грани ($\chi=1$ для горизонтальных граней, $\chi=0{,}5$ для боковых); $\mathcal T_R=\exp(-\tau_R m_a)$ --- пропускание Рэлея, $\tau_R=0{,}09\,msl/101325$ --- оптическая толщина (на уровне моря $0{,}09$), $m_a=\min(1/\cos Z,\,40)$ --- воздушная масса; $\mathcal T_{wv}=1-0{,}077\,w^{0{,}3}$ --- пропускание водяного пара ($w$ --- влагосодержание столба, пропорциональное $e_{v}(101325/msl)$; $e_v$ --- парциальное давление водяного пара, вычисленное по $t_{2m}$ и $d_{2m}$ через насыщение по формуле Тетенса); $\mathcal T_a=0{,}93$ --- аэрозоли (арктический фон); $\mathcal T_c=1-0{,}75\,tcc$ --- облачность.

\paragraph{Длинноволновая радиация.}
\begin{equation}
L_\downarrow=\sigma_a T_a^4\,(1+0{,}275\,tcc)\Bigl[1-0{,}261\,e^{-7{,}77\cdot10^{-4}(273{,}15-T_a)^2}\Bigr],\quad\sigma_a=5{,}4999\cdot10^{-8}\ \text{Вт/(м}^2\text{К}^4),
\end{equation}
\begin{equation}
Q_{LW}=\varepsilon\sigma\left(T_a^4-T_s^4\right)+(1-\varepsilon)L_\downarrow,\qquad
\frac{\partial Q_{LW}}{\partial T_s}=-4\varepsilon\sigma T_s^3,
\end{equation}
где $T_a=t_{2m}$, К; $T_s$ --- температура поверхности, К; $\varepsilon=0{,}97$; $\sigma=5{,}67\cdot10^{-8}$~Вт/(м\textsuperscript2К\textsuperscript4).

\paragraph{Явный и скрытый потоки.}
\begin{equation}
Q_{SH}=\rho_a c_{pa}C_H\,U\,(T_a-T_s),\qquad
Q_{LH}=\rho_a L_s C_E\,U\,\bigl(q_a-q_{\mathrm{sat}}(T_s)\bigr),
\end{equation}
$c_{pa}=1004$~Дж/(кг$\cdot$К), $L_s=2{,}835\cdot10^6$~Дж/кг, $C_H=C_E=1{,}5\cdot10^{-3}$ (фиксированные нейтральные коэффициенты), $U$ --- скорость ветра на 10~м с нижним порогом, $q_a=0{,}622\,e_v/msl$ --- удельная влажность воздуха, $\rho_a=msl/(287\,T_a)$. Насыщающая упругость над льдом
\begin{equation}
e_{\mathrm{sat}}=610{,}78\exp\!\left(\frac{21{,}875\,(T_s-273{,}15)}{T_s-7{,}66}\right),\qquad
q_{\mathrm{sat}}=\frac{0{,}622\,e_{\mathrm{sat}}}{p}.
\end{equation}
Производные $\partial Q_{SH}/\partial T_s=-\rho_ac_{pa}C_HU$ и $\partial Q_{LH}/\partial T_s$ (по конечной разности $q_{\mathrm{sat}}$ на $\pm0{,}5$~К) входят в диагональ тепловой матрицы. Масса, теряемая при сублимации, не учитывается: $Q_{LH}$ --- только энергетический поток.

\paragraph{Что не учтено.} Проникновение коротковолновой радиации в лёд (вся поглощённая радиация ставится в поверхностный воксель); затенение одних граней другими; зависимость от наклона и азимута грани (вместо неё --- $\chi=1$ либо $0{,}5$); изменение альбедо во времени (мокрый лёд, талые лужи); дождь и снег как источники тепла и массы; суточный ход (используются суточные средние ERA5).

%=====================================================================
\section{Фазовый переход и ограничение таяния}\label{sec:melt}

\subsection{Фазовое ограничение}
Если после решения тепловой задачи температура вокселя положительна, избыток энтальпии превращается в таяние:
\begin{equation}
\Delta\varphi_i=\min\!\left(\varphi_i\frac{\ci\,T_i}{\Lf},\,\varphi_i\right),\qquad T_i\leftarrow0,
\end{equation}
масса растаявшего $\rhoi\Delta\varphi_iV$ учитывается как поверхностное таяние. Множитель $\varphi_i$ делает преобразование энергетически согласованным с $C_i=\varphi_i\rhoi\ci V$.

\subsection{Граничное таяние и ограничение за шаг}
Накопленный по граням объём таяния $\Delta\varphi_i^{\mathrm{bnd}}$ применяется после тепловой задачи:
\begin{equation}
\varphi_i\leftarrow\varphi_i-\min\bigl(\Delta\varphi_i^{\mathrm{bnd}},\,\mu\,\varphi_i\bigr),\qquad\mu=0{,}20 ,
\end{equation}
где $\mu$ --- предельная доля таяния вокселя за шаг (\code{melt\_cap}); при срабатывании ограничения базальная и боковая части масштабируются пропорционально. Для $h=1$~м и $\Delta t=3600$~с это соответствует предельной скорости отступания грани $\approx5{,}6\cdot10^{-5}$~м/с ($\approx4{,}8$~м/сут).

\subsection{Слив и уплотнение}
Воксели с $\varphi<\varphi_{\min}$ сливаются (раздел~\ref{sec:flush}); после уплотнения вызывается пересчёт гидростатики.

%=====================================================================
\section{Баланс энергии и массы}\label{sec:ledger}

Для контроля ведётся бухгалтерия энергии: интеграл граничного теплового потока за шаг по типам граней с тем же $\theta$-взвешиванием, что и в схеме, $E=\Delta t\,[\theta q^{n+1}+(1-\theta)q^{n}]$. Остаток
\begin{equation}
R=Q_{\mathrm{bnd}}-Q_{\mathrm{melt,sens}}-\bigl[Q_{\mathrm{calved}}+Q_{\mathrm{flushed}}+\Lf\,m_{\mathrm{surf}}\bigr],
\end{equation}
где $Q_{\mathrm{bnd}}$ --- суммарный граничный поток; $Q_{\mathrm{melt,sens}}$ --- сенсибельное тепло льда, ушедшее с граничным таянием; $Q_{\mathrm{calved}}$, $Q_{\mathrm{flushed}}$ --- энергия откольных и слитых вокселей; $m_{\mathrm{surf}}$ --- масса, растаявшая по фазовому ограничению. Относительный остаток на проверенных расчётах: $0{,}0166$--$0{,}0195$~\%.

%=====================================================================
\section{Грунтование и условия остановки}

Айсберг считается севшим на мель, если осадка превышает глубину IBCAO в данной точке. В приложении \code{iceberg\_run.f90} предусмотрены режимы: \code{--strict-ocean} (остановка при выходе в точку без океанских данных), \code{--stop-on-grounding} (остановка при грунтовании), \code{--grounded-in-place} (айсберг остаётся на месте и продолжает таять). Вследствие остановок у пяти из шести прибрежных прогонов срок жизни является оценкой снизу; лишь прогон FJL\_A\_2020 растаял полностью (323{,}8 сут).

%=====================================================================
\section{Волны}\label{sec:waves}

\textbf{Текущее состояние.} Волновая эрозия в формулировке Бигга и др. \cite{bigg1997}, принятой в NEMO-ICB \cite{marsh2015},
\begin{equation}
m_{\mathrm{wave}}=\frac{1}{12}\,(T_w+2)\,S_s\,\bigl[1+\cos(\pi c^3)\bigr]\ \ \text{м/сут},
\end{equation}
где $T_w$ --- температура поверхности воды, \textdegree C, $c$ --- сплочённость морского льда, $S_s$ --- состояние моря (по шкале Дугласа из высоты значительных волн $H_s$), реализована в коде, но \textbf{выключена} по умолчанию (\code{ICEBERG\_WAVE\_EROSION}, \code{ICEBERG\_STOKES\_DRIFT}). Причины: отсутствие согласованных полей волн (имеющиеся данные WAM несовместимы по формату с \code{wave\_forcing.f90}), а также то, что эрозия реализована на уровне призматического представления (уменьшение длины и ширины), а не на воксельной геометрии. В текущих воксельных расчётах волновое таяние и волновой дрейф \emph{отсутствуют}.

\textbf{Как волны вписываются в схему интерфейса} (проект развития, раздел~\ref{sec:future}): через замену $u_{\mathrm{rel}}\to\sqrt{u_{\mathrm{rel}}^2+u_{\mathrm{orb}}^2}$ в $\gamma_T,\gamma_S$ и через дополнительную эрозию вокселей ватерлинии.

%=====================================================================
\section{Верификация}

\begin{itemize}
\item Сходимость по сетке: на $100^3$ ошибка $3{,}75$~\%, оценённый порядок $p=0{,}711$.
\item Ретроспективный расчёт по данным Эндерлин и др. \cite{enderlin2023}: 190 из 190 случаев; медиана отношения модель/наблюдение $0{,}540$, ранговая корреляция Спирмена $0{,}156$ (сетка $8^3$). Модель в среднем занижает таяние и слабо ранжирует отдельные айсберги.
\item Остаток энергии $0{,}0166$--$0{,}0195$~\% (раздел~\ref{sec:ledger}).
\item Дрейф по наблюдениям не валидирован.
\end{itemize}

%=====================================================================
\section{Упрощения и заглушки}\label{sec:stubs}

\begin{longtable}{@{}p{3.1cm}p{6.2cm}p{5.6cm}@{}}
\toprule
Процесс & Как в модели сейчас & Последствие\\
\midrule
\endhead
Форма айсберга & начальный куб $100^3$~м; перетяжки и «грибообразные» свесы принимаются как есть (KNOWN\_ISSUES Q-31, D-81, D-82) & площади боковых граней и киля завышены/занижены относительно реальных форм; прямое влияние на дрейф и таяние\\
Опрокидывание & только оценка $GM$, ориентация не меняется & нет резкого изменения экспозиции граней\\
Волны & выключены (раздел~\ref{sec:waves}) & отсутствует волновое таяние ватерлинии, волновой дрейф\\
Морской лёд & нет сопротивления, нет экранирования & переоценка дрейфа в ледовых условиях\\
Радиация & суточные средние, без проникновения, затенения, наклона & нет суточного хода и объёмного нагрева\\
Фрагменты отколов & не прослеживаются & судьба откола неизвестна\\
Градиент давления & нулевой (течения реанализа) & геострофика только в течениях\\
Аккреция & $m\ge0$ & нет намерзания\\
Талая вода & не возвращается в океан & нет обратной связи по $T_w,S_w$ у борта\\
Константы & не калиброваны; часть без привязки к первоисточнику (параметры атмосферной прозрачности, формула $e_{\mathrm{sat}}$, свободная конвекция) & нужен анализ чувствительности\\
Дрейф & не валидирован & нужны сопоставления с треками\\
Начальные данные & стартовые точки приблизительные & нужны точки по наблюдениям\\
Грунтование и суша & остановка или оставление на месте & нижние оценки сроков жизни\\
Разрешение & $h=1$~м, не разрешает суточный ход температуры поверхности & занижение суточных экстремумов потока\\
\bottomrule
\end{longtable}

%=====================================================================
\section{Развитие модели}\label{sec:future}

Раздел содержит проекты, в коде они не реализованы. Для каждого указано, где именно изменение входит в текущие уравнения.

\subsection{Геометрия айсберга}
\begin{enumerate}
\item \textbf{Реалистичные начальные формы.} Начальные поля $\varphi_i$ и $T_i$ уже являются произвольными массивами: форма задаётся маской вокселей. Достаточно построить генераторы форм (таблитчатая с отношением длины к высоте по наблюдениям, клиновидная, пирамидальная, «купольная») и подавать их вместо куба; при сохранении $\Delta t$, схемы CN и интерфейса никаких изменений уравнений не требуется. Данные о форме и деградации ледяных островов Канадской Арктики приведены, в частности, в работе \cite{crawford2013}.
\item \textbf{Подъячеечная геометрия.} Значение $\varphi\in(0,1)$ можно трактовать как заполнение по объёму и вычислять площадь граничной грани через оценку нормали по градиенту $\nabla\varphi$ (метод «объёма жидкости»); это уменьшает ступенчатость при малом $h$ и порядке сходимости $p<2$.
\item \textbf{Опрокидывание.} При $GM<0$ или нарушении критерия устойчивости \cite{wagner2017} тело поворачивается (перенос $\varphi$ и $T$ на повёрнутую сетку или смена ориентации матрицы граней); входит как дискретное событие, после которого включается процедура Раннахера.
\item \textbf{Неравномерная и вложенная сетка.} Для крупных тел $h=1$~м избыточен в глубине и недостаточен у поверхности; возможны вложенные слои у ватерлинии и поверхности.
\item \textbf{Фрагменты откола} как отдельные лагранжевы частицы с собственной массой и энергией (сейчас учитываются только в бухгалтерии).
\end{enumerate}

\subsection{Солнечная радиация и воздушный баланс}
\begin{enumerate}
\item \textbf{Проникновение коротковолновой радиации.} Заменить поверхностный источник объёмным по закону Бугера--Ламберта: в правую часть вокселя добавляется
\[
Q_{i}^{\mathrm{SW}}=(1-\alpha)\,I_0\,\bigl(e^{-\kappa s_i}-e^{-\kappa (s_i+h)}\bigr)A,
\]
где $I_0$ --- поглощённая у поверхности радиация, $\kappa$ --- коэффициент ослабления в льду, $s_i$ --- глубина верхней грани вокселя вдоль луча, $A$ --- площадь грани. Член входит в $r_i$ и не меняет матрицу.
\item \textbf{Затенение и ориентация.} Вместо $\chi$ использовать $\chi_f=\max(0,\mathbf n_f\!\cdot\mathbf s)$ ($\mathbf n_f$ --- нормаль грани, $\mathbf s$ --- направление на Солнце) и отслеживать луч до первого заполненного вокселя на маске $\varphi$.
\item \textbf{Использовать SSRD и STRD ERA5} вместо параметризации (при наличии почасовых данных).
\item \textbf{Меняющееся альбедо} (мокрый лёд, талые лужи) и учёт осадков.
\end{enumerate}

\subsection{Волны}
\begin{enumerate}
\item \textbf{Волновая конвекция в интерфейсе.} Орбитальная скорость $u_{\mathrm{orb}}(z)=\dfrac{\pi H_s}{T_p}\,e^{-kz}$ ($H_s$, $T_p$ --- высота и период волн, $k$ --- волновое число) добавляется в относительную скорость: $u_{\mathrm{rel}}^{\mathrm{eff}}=\sqrt{u_{\mathrm{rel}}^2+u_{\mathrm{orb}}^2}$; уравнения (\ref{eq:salt})--(\ref{eq:heat}) не меняются, увеличиваются $\gamma_T$, $\gamma_S$ у ватерлинии.
\item \textbf{Механическая эрозия ватерлинии.} Скорость $m_{\mathrm{wave}}$ применяется к граням воксельного слоя у $z_{wl}$ как дополнительный вклад в $\Delta\varphi_i^{\mathrm{bnd}}$ с тем же ограничением $\mu$.
\item \textbf{Стоксов дрейф} и волновое радиационное напряжение входят в $\mathbf F_w$ как добавка к течению $\mathbf u_k$ верхних слоёв.
\item Поля: ERA5-волны или WAM, приведённые к формату \code{wave\_forcing.f90}.
\end{enumerate}

\subsection{Интерфейс и теплообмен}
\begin{enumerate}
\item \textbf{Аккреция.} Допустить $m<0$ при $F(0)<0$ (намерзание на холодных гранях), введя баланс для солёности и теплоты кристаллизации.
\item \textbf{Нелинейная зависимость $\gamma$ от $u_{\mathrm{rel}}$ и шероховатости} (формулы Нуссельта) вместо $K_Tu_{\mathrm{rel}}$.
\item \textbf{Обратная связь по талой воде:} вблизи борта $T_w,S_w$ заменяются смесью океанской воды и талой воды (одномерная пластина)
\end{enumerate}

\subsection{Прочее}
Калибровка и анализ чувствительности ($K_T$, $K_S$, $\gamma^{\mathrm{free}}$, параметры радиации); валидация дрейфа по трекам; сценарные эксперименты (сдвиг температуры океана на $\pm2$--$3$~\textdegree C относительно 2020--2024); режимы обхода мелководья.

%=====================================================================
\section{Основные параметры и константы}

\begin{longtable}{@{}p{5.8cm}p{3.6cm}p{5.5cm}@{}}
\toprule
Параметр & Значение & Примечание\\
\midrule
\endhead
Сетка, шаг $h$ & $100^3$, 1~м & базовая конфигурация\\
$\varphi_{\min}$, $\varepsilon_\varphi$ & $0{,}01$; $10^{-12}$ & порог слива; активность\\
$T_0$ & $-10$~\textdegree C & начальная температура\\
$\Delta t$ & 3600~с & тепло и дрейф\\
$\theta$ & 0{,}5 & Кранк--Николсон\\
Процедура Раннахера & 2 шага $\theta=1$ & после смены геометрии\\
Допуск PCG, макс. итераций & $10^{-10}$; $10\,N_xN_yN_z$ (env override) & Якоби; \code{ICEBERG\_CG\_TOL}, \code{ICEBERG\_CG\_MAXIT}\\
$\mu$ (\code{melt\_cap}) & 0{,}20 & доля $\varphi$ за шаг\\
Порог связности, период & $\varphi>0{,}3$; по умолчанию 1~ч, максимум 6~ч; в 12.20 --- 6~ч & \code{VOXEL\_PHI\_CONN}, \code{VOXEL\_CONNECTIVITY\_HOURS}\\
$\rhoi$, $\rhow$, $\rho_a$ & 910; 1028; 1{,}225 кг/м\textsuperscript3 & плотности\\
$\ci$, $c_w$, $c_{pa}$ & 2100; 3974; 1004 Дж/(кг$\cdot$К) & теплоёмкости\\
$\ki$, $\Lf$, $L_s$ & 2{,}2 Вт/(м$\cdot$К); $3{,}34\cdot10^5$; $2{,}835\cdot10^6$ Дж/кг & \\
$K_T$, $K_S$ & $1{,}1\cdot10^{-3}$; $3{,}1\cdot10^{-5}$ & \cite{jenkins2010}\\
$\gamma^{\mathrm{free}}_T$, $\gamma^{\mathrm{free}}_S$ & $10^{-5}$; $10^{-6}$ м/с & нижние пределы\\
$C_a$, $C_w$ & 1{,}3; 0{,}9 & форменное сопротивление\\
Кожное: воздух; вода & $2{,}2\cdot10^{-3}$; $5{,}5\cdot10^{-3}$ & \\
$\alpha$, $\varepsilon$ & 0{,}7; 0{,}97 & альбедо; излучательная способность\\
$C_H=C_E$ & $1{,}5\cdot10^{-3}$ & нейтральные коэффициенты\\
$\chi$ & 1; 0{,}5 & горизонтальные; боковые грани\\
\bottomrule
\end{longtable}

%=====================================================================
\begin{thebibliography}{99}
\bibitem{bigg1997} Bigg G.\,R., Wadley M.\,R., Stevens D.\,P., Johnson J.\,A. Modelling the dynamics and thermodynamics of icebergs // Cold Regions Science and Technology. 1997. Vol.~26, no.~2. P.~113--135. DOI: 10.1016/S0165-232X(97)00012-8.
\bibitem{keghouche2009} Keghouche I., Bertino L., Lisæter K.\,A. Parameterization of an iceberg drift model in the Barents Sea // Journal of Atmospheric and Oceanic Technology. 2009. Vol.~26, no.~10. P.~2216--2227. DOI: 10.1175/2009JTECHO678.1.
\bibitem{marsh2015} Marsh R., Ivchenko V., Skliris N. et al. NEMO--ICB (v1.0): interactive icebergs in the NEMO ocean model globally configured at eddy-permitting resolution // Geoscientific Model Development. 2015. Vol.~8. P.~1547--1562. DOI: 10.5194/gmd-8-1547-2015.
\bibitem{crawford2013} Crawford A. Ice island deterioration in the Canadian Arctic: rates, patterns and model evaluation. M.\,Sc. thesis. Carleton University, 2013. DOI: 10.22215/etd/2013-10024.
\bibitem{wagner2017} Wagner T.\,J.\,W., Stern A.\,A., Dell R.\,W., Eisenman I. On the representation of capsizing in iceberg models // Ocean Modelling. 2017. Vol.~117. P.~88--96. DOI: 10.1016/j.ocemod.2017.07.003.
\bibitem{enderlin2023} Enderlin E.\,M., Moffat C., Miller E. et al. Antarctic iceberg melt rate variability and sensitivity to ocean thermal forcing // Journal of Glaciology. 2023. Vol.~69, no.~278. P.~1718--1728. DOI: 10.1017/jog.2023.54.
\bibitem{hollandjenkins1999} Holland D.\,M., Jenkins A. Modeling thermodynamic ice--ocean interactions at the base of an ice shelf // Journal of Physical Oceanography. 1999. Vol.~29, no.~8. P.~1787--1800. DOI: 10.1175/1520-0485(1999)029<1787:MTIOIA>2.0.CO;2.
\bibitem{jenkins2010} Jenkins A., Nicholls K.\,W., Corr H.\,F.\,J. Observation and parameterization of ablation at the base of Ronne Ice Shelf, Antarctica // Journal of Physical Oceanography. 2010. Vol.~40, no.~10. P.~2298--2312. DOI: 10.1175/2010JPO4317.1.
\bibitem{hellmer1989} Hellmer H.\,H., Olbers D.\,J. A two-dimensional model for the thermohaline circulation under an ice shelf // Antarctic Science. 1989. Vol.~1, no.~4. P.~325--336.
\bibitem{spencer1971} Spencer J.\,W. Fourier series representation of the position of the Sun // Search. 1971. Vol.~2, no.~5. P.~172.
\bibitem{crank1947} Crank J., Nicolson P. A practical method for numerical evaluation of solutions of partial differential equations of the heat-conduction type // Proceedings of the Cambridge Philosophical Society. 1947. Vol.~43, no.~1. P.~50--67.
\bibitem{rannacher1984} Rannacher R. Finite element solution of diffusion problems with irregular data // Numerische Mathematik. 1984. Vol.~43. P.~309--327.
\bibitem{hestenes1952} Hestenes M.\,R., Stiefel E. Methods of conjugate gradients for solving linear systems // Journal of Research of the National Bureau of Standards. 1952. Vol.~49, no.~6. P.~409--436.
\bibitem{samarskii1977} Самарский А.\,А. Теория разностных схем. М.: Наука, 1977. 656 с.
\end{thebibliography}

\end{document}
